\documentclass[11pt,a4paper,fontset=mac]{ctexart}
\usepackage{geometry}\geometry{margin=2.4cm}
\usepackage{amsmath,booktabs,threeparttable,setspace}
\usepackage{xcolor}
\usepackage[colorlinks=true,linkcolor=blue!60!black]{hyperref}
\onehalfspacing
\title{\textbf{时序可择时性的来源:跨资产 AR--GARCH 证据}\\
\large ——中证1000 择时策略为何有效,以及在何处可以复制}
\author{}
\date{2026年7月}

\begin{document}
\maketitle
\vspace{-2.5em}
\begin{abstract}
\noindent 本文以 23 个资产(A股宽基指数、贵金属、商品期货、海外指数与 A 股个股)的日频数据估计带状态交互项的 AR(1)--GARCH(1,1) 模型,考察时序择时能力的统计来源。主要发现:(i) 短期动量系数 $\varphi_1$ 与实测择时策略绩效的秩相关达 $0.77$($p<0.01$),是可择时性的第一判别变量;(ii) 条件波动的水平项 $\varphi_2$(``波动溢价'')在中证1000 上约为 $+11$bp/日(每标准差),为全样本最强,与遗传规划因子库中波动族因子占比约六成的事实互为印证;(iii) 波动可预测性(GARCH 持续性、Mincer--Zarnowitz $R^2$)与择时绩效\emph{无关};(iv) 个股层面 $\varphi_1$ 一致为负,指数层面的正自相关是横截面聚合(lead--lag)的涌现性质;(v) 样本外检验:国证2000 与中证2000 呈现与中证1000 几乎相同且更强的参数结构,是本框架最优先的推广标的。
\end{abstract}

\section{模型设定}
对数收益 $r_t$(单位:\%)的均值方程与波动方程为
\begin{align}
r_t &= \mu + \varphi_1 r_{t-1} + \varphi_L\, \tilde m_{t-1} + \varphi_2\, \tilde\sigma_{t-1} + \varphi_3\,(\tilde\sigma_{t-1}\cdot r_{t-1}) + \varepsilon_t, \\
\varepsilon_t &= \sigma_t e_t,\qquad e_t \sim \mathrm{i.i.d.}\ N(0,1),\qquad
\sigma_t^2 = \omega + \alpha\,\varepsilon_{t-1}^2 + \beta\,\sigma_{t-1}^2,
\end{align}
其中 $\tilde\sigma_{t-1}$ 为条件波动的标准化值(状态变量),$\tilde m_{t-1}$ 为过去 244 个交易日累计收益的标准化值(长期动量),交互项 $\tilde\sigma_{t-1}\cdot r_{t-1}$ 度量波动状态对短期自相关的调制:$\varphi_1 + \varphi_3\tilde\sigma_{t-1}$ 即时变一阶自相关。

估计采用两步法:第一步以 AR(1)--GARCH(1,1) 提取条件波动 $\hat\sigma_t$ 并构造状态变量;第二步将其作为外生回归量与交互项一并进入均值方程,扰动仍设 GARCH(1,1),以 QMLE 估计,报告 Bollerslev--Wooldridge 稳健 $t$ 值。全部序列截取 2010 年以后,剔除 $|r_t|>25\%$ 的异常观测。

\paragraph{解读约定}$\varphi_1>0$:日频动量(昨日上涨预示今日续涨);$\varphi_2>0$:高波动状态具有正的风险补偿(``恐慌后反弹'');$\varphi_L>0$:年度动量的延续;$\varphi_3<0$:高波动状态下反转增强。

\section{数据}
A股指数与商品主力连续、黄金现货(au9999)来自本项目行情库;国证2000、中证2000、北证50 与六只 A 股个股(前复权)来自东方财富历史行情接口。个股按市值与投资者结构选取谱系:贵州茅台(超大盘)、中国平安(大盘)、华鲁恒升(中盘)、中际旭创(创业板大盘)、利欧股份与光华科技(小盘)。

\section{估计结果}
\begin{table}[htbp]\centering
\caption{A股宽基指数:五项均值方程估计(GARCH(1,1) 扰动)}\label{tab:ashare}
\footnotesize\setlength{\tabcolsep}{4pt}
\begin{threeparttable}
\begin{tabular}{lccccccc}
\toprule
 & 沪深300 & 中证500 & 中证1000 & 创业板指 & 国证2000 & 中证2000 & 北证50 \\
\midrule
$\hat\mu$~常数项 (年化\%) & $3.0$ & $6.5$ & $10.1$ & $13.9^{*}$ & $13.0^{*}$ & $11.3$ & $-4.3$ \\
 & $(0.55)$ & $(1.03)$ & $(1.46)$ & $(1.81)$ & $(1.92)$ & $(1.42)$ & $(-0.17)$ \\[2pt]
$\hat\varphi_1$~($r_{t-1}$) & $0.0151$ & $0.0376^{*}$ & $0.0550^{***}$ & $0.0312^{*}$ & $0.0538^{***}$ & $0.0540^{**}$ & $-0.0051$ \\
 & $(0.84)$ & $(1.95)$ & $(2.70)$ & $(1.72)$ & $(2.61)$ & $(2.32)$ & $(-0.10)$ \\[2pt]
$\hat\varphi_L$~($\tilde m_{t-1}$, bp) & $1.2$ & $5.2$ & $5.4$ & $4.9$ & $7.5^{**}$ & $10.6^{**}$ & $23.8$ \\
 & $(0.44)$ & $(1.49)$ & $(1.47)$ & $(1.08)$ & $(2.11)$ & $(2.53)$ & $(0.81)$ \\[2pt]
$\hat\varphi_2$~($\tilde\sigma_{t-1}$, bp) & $-1.4$ & $6.0^{*}$ & $11.0^{***}$ & $7.9^{**}$ & $10.8^{***}$ & $10.8^{**}$ & $10.8$ \\
 & $(-0.52)$ & $(1.81)$ & $(2.84)$ & $(2.03)$ & $(2.92)$ & $(2.46)$ & $(0.83)$ \\[2pt]
$\hat\varphi_3$~($\tilde\sigma_{t-1}\!\cdot\! r_{t-1}$) & $-0.0206$ & $0.0032$ & $0.0051$ & $-0.0035$ & $0.0019$ & $0.0009$ & $0.0092$ \\
 & $(-1.17)$ & $(0.17)$ & $(0.25)$ & $(-0.21)$ & $(0.09)$ & $(0.04)$ & $(0.29)$ \\[2pt]
\midrule
$\hat\alpha+\hat\beta$ & $0.981$ & $0.952$ & $0.933$ & $0.978$ & $0.929$ & $0.909$ & $0.975$ \\
MZ $R^2$ & $0.115$ & $0.138$ & $0.135$ & $0.106$ & $0.114$ & $0.112$ & $0.110$ \\
VR(5) & $1.022$ & $1.100$ & $1.147$ & $1.059$ & $1.156$ & $1.205$ & $1.113$ \\
买持夏普 & $0.07$ & $0.13$ & $0.20$ & $0.26$ & $0.20$ & $0.30$ & $0.05$ \\
年化波动(\%) & $21$ & $25$ & $27$ & $30$ & $26$ & $28$ & $38$ \\
$N$ & $3771$ & $3771$ & $3527$ & $3671$ & $3530$ & $2806$ & $853$ \\
样本起点 & \scriptsize 2010-01 & \scriptsize 2010-01 & \scriptsize 2010-01 & \scriptsize 2010-06 & \scriptsize 2011-01 & \scriptsize 2014-01 & \scriptsize 2022-05 \\
\bottomrule
\end{tabular}
\begin{tablenotes}\scriptsize
\item 括号内为 Bollerslev--Wooldridge 稳健 $t$ 值;$^{***}$、$^{**}$、$^{*}$ 分别表示 1\%、5\%、10\% 显著水平(双侧)。
\item $\hat\varphi_L,\hat\varphi_2$ 单位为 bp/日(状态变量每高一个标准差);长动量窗为 $\min(244,N/6)$ 日。
\end{tablenotes}\end{threeparttable}\end{table}

\begin{table}[htbp]\centering
\caption{贵金属与商品期货(主力连续)}\label{tab:comm}
\footnotesize\setlength{\tabcolsep}{4pt}
\begin{threeparttable}
\begin{tabular}{lcccccccc}
\toprule
 & 黄金现货 & 黄金期货 & 白银 & 螺纹钢 & 铜 & 豆粕 & PTA & 甲醇 \\
\midrule
$\hat\mu$~常数项 (年化\%) & $15.0^{***}$ & $6.7^{*}$ & $6.6$ & $-1.5$ & $4.7$ & $-2.9$ & $-3.2$ & $4.7$ \\
 & $(3.27)$ & $(1.70)$ & $(0.90)$ & $(-0.26)$ & $(1.04)$ & $(-0.45)$ & $(-0.54)$ & $(0.55)$ \\[2pt]
$\hat\varphi_1$~($r_{t-1}$) & $0.0156$ & $-0.0079$ & $0.0058$ & $-0.0440^{**}$ & $-0.0438^{**}$ & $-0.0530^{**}$ & $0.0130$ & $-0.0436^{**}$ \\
 & $(0.62)$ & $(-0.42)$ & $(0.29)$ & $(-2.32)$ & $(-2.53)$ & $(-2.12)$ & $(0.68)$ & $(-2.10)$ \\[2pt]
$\hat\varphi_L$~($\tilde m_{t-1}$, bp) & $0.1$ & $1.5$ & $-1.5$ & $1.1$ & $0.5$ & $-5.4$ & $-0.8$ & $-0.1$ \\
 & $(0.04)$ & $(0.82)$ & $(-0.39)$ & $(0.34)$ & $(0.26)$ & $(-1.57)$ & $(-0.32)$ & $(-0.03)$ \\[2pt]
$\hat\varphi_2$~($\tilde\sigma_{t-1}$, bp) & $2.0$ & $-0.5$ & $2.4$ & $2.2$ & $1.2$ & $-0.9$ & $2.4$ & $-2.7$ \\
 & $(0.46)$ & $(-0.28)$ & $(0.56)$ & $(0.82)$ & $(0.55)$ & $(-0.30)$ & $(1.02)$ & $(-0.69)$ \\[2pt]
$\hat\varphi_3$~($\tilde\sigma_{t-1}\!\cdot\! r_{t-1}$) & $-0.0216$ & $-0.0188$ & $-0.0111$ & $-0.0251$ & $-0.0473^{***}$ & $0.0236$ & $-0.0659^{***}$ & $-0.0253$ \\
 & $(-0.87)$ & $(-0.97)$ & $(-0.61)$ & $(-1.62)$ & $(-3.03)$ & $(1.31)$ & $(-3.97)$ & $(-1.45)$ \\[2pt]
\midrule
$\hat\alpha+\hat\beta$ & $1.046$ & $1.006$ & $1.007$ & $0.987$ & $0.951$ & $1.015$ & $1.001$ & $1.006$ \\
MZ $R^2$ & $0.082$ & $0.072$ & $0.185$ & $0.087$ & $0.104$ & $0.009$ & $0.079$ & $0.100$ \\
VR(5) & $0.930$ & $0.922$ & $1.058$ & $0.935$ & $0.952$ & $0.882$ & $0.975$ & $0.916$ \\
买持夏普 & $0.84$ & $0.48$ & $0.21$ & $-0.10$ & $0.18$ & $0.00$ & $-0.09$ & $0.11$ \\
年化波动(\%) & $15$ & $16$ & $28$ & $23$ & $19$ & $22$ & $24$ & $29$ \\
$N$ & $2077$ & $3776$ & $3206$ & $3767$ & $3773$ & $3771$ & $3764$ & $2556$ \\
样本起点 & \scriptsize 2016-12 & \scriptsize 2010-01 & \scriptsize 2012-05 & \scriptsize 2010-01 & \scriptsize 2010-01 & \scriptsize 2010-01 & \scriptsize 2010-01 & \scriptsize 2014-12 \\
\bottomrule
\end{tabular}
\begin{tablenotes}\scriptsize
\item 括号内为 Bollerslev--Wooldridge 稳健 $t$ 值;$^{***}$、$^{**}$、$^{*}$ 分别表示 1\%、5\%、10\% 显著水平(双侧)。
\item $\hat\varphi_L,\hat\varphi_2$ 单位为 bp/日(状态变量每高一个标准差);长动量窗为 $\min(244,N/6)$ 日。
\end{tablenotes}\end{threeparttable}\end{table}

\begin{table}[htbp]\centering
\caption{海外指数与 A 股个股(前复权)}\label{tab:stk}
\footnotesize\setlength{\tabcolsep}{4pt}
\begin{threeparttable}
\begin{tabular}{lcccccccc}
\toprule
 & 纳斯达克100 & 标普500 & 贵州茅台 & 中国平安 & 华鲁恒升 & 利欧股份 & 中际旭创 & 光华科技 \\
\midrule
$\hat\mu$~常数项 (年化\%) & $21.8^{***}$ & $18.7^{***}$ & $39.6^{**}$ & $20.8$ & $74.0^{***}$ & $-4.9$ & $67.8^{***}$ & $-4.4$ \\
 & $(3.65)$ & $(4.72)$ & $(2.20)$ & $(0.97)$ & $(2.89)$ & $(-0.37)$ & $(2.65)$ & $(-0.28)$ \\[2pt]
$\hat\varphi_1$~($r_{t-1}$) & $-0.0309$ & $-0.0403^{**}$ & $-0.0193$ & $-0.0236$ & $-0.0552^{***}$ & $0.0206$ & $0.0629^{***}$ & $0.0097$ \\
 & $(-1.39)$ & $(-1.96)$ & $(-0.79)$ & $(-1.04)$ & $(-2.73)$ & $(1.02)$ & $(2.88)$ & $(0.42)$ \\[2pt]
$\hat\varphi_L$~($\tilde m_{t-1}$, bp) & $-1.0$ & $-1.9$ & $-2.5$ & $-4.0$ & $-11.8$ & $6.4$ & $6.2$ & $-2.9$ \\
 & $(-0.39)$ & $(-1.39)$ & $(-0.26)$ & $(-0.57)$ & $(-0.69)$ & $(0.97)$ & $(0.73)$ & $(-0.45)$ \\[2pt]
$\hat\varphi_2$~($\tilde\sigma_{t-1}$, bp) & $5.4$ & $6.7^{**}$ & $25.7^{*}$ & $10.8$ & $44.8^{**}$ & $1.5$ & $20.0$ & $-6.3$ \\
 & $(1.62)$ & $(2.53)$ & $(1.73)$ & $(0.67)$ & $(1.97)$ & $(0.19)$ & $(1.18)$ & $(-0.69)$ \\[2pt]
$\hat\varphi_3$~($\tilde\sigma_{t-1}\!\cdot\! r_{t-1}$) & $-0.0605^{***}$ & $-0.0322^{**}$ & $0.0132$ & $-0.0257$ & $-0.0409^{*}$ & $0.0469^{**}$ & $0.0318$ & $0.0075$ \\
 & $(-3.90)$ & $(-2.22)$ & $(0.32)$ & $(-0.69)$ & $(-1.86)$ & $(2.47)$ & $(1.16)$ & $(0.29)$ \\[2pt]
\midrule
$\hat\alpha+\hat\beta$ & $0.890$ & $0.836$ & $0.985$ & $0.987$ & $1.001$ & $0.975$ & $0.993$ & $0.923$ \\
MZ $R^2$ & $0.267$ & $0.352$ & $0.100$ & $0.186$ & $0.299$ & $0.240$ & $0.203$ & $0.154$ \\
VR(5) & $0.826$ & $0.857$ & $0.947$ & $0.993$ & $0.837$ & $1.204$ & $1.081$ & $1.071$ \\
买持夏普 & $0.76$ & $0.66$ & $0.45$ & $0.30$ & $0.59$ & $0.21$ & $0.77$ & $0.25$ \\
年化波动(\%) & $21$ & $17$ & $53$ & $69$ & $86$ & $53$ & $83$ & $52$ \\
$N$ & $2875$ & $3917$ & $2161$ & $2817$ & $2775$ & $3435$ & $2726$ & $2485$ \\
样本起点 & \scriptsize 2014-02 & \scriptsize 2010-01 & \scriptsize 2016-07 & \scriptsize 2010-01 & \scriptsize 2010-01 & \scriptsize 2010-01 & \scriptsize 2012-04 & \scriptsize 2015-02 \\
\bottomrule
\end{tabular}
\begin{tablenotes}\scriptsize
\item 括号内为 Bollerslev--Wooldridge 稳健 $t$ 值;$^{***}$、$^{**}$、$^{*}$ 分别表示 1\%、5\%、10\% 显著水平(双侧)。
\item $\hat\varphi_L,\hat\varphi_2$ 单位为 bp/日(状态变量每高一个标准差);长动量窗为 $\min(244,N/6)$ 日。
\end{tablenotes}\end{threeparttable}\end{table}

\begin{table}[htbp]\centering
\caption{模型参数与实测策略效果的 Spearman 秩相关($n=12$)}\label{tab:rank}
\small\begin{threeparttable}\begin{tabular}{lcc}\toprule
参数 & $\hat\rho$ & $p$ 值 \\ \midrule
$\hat\varphi_1$~($r_{t-1}$,短动量) & $+0.766$$^{***}$ & $0.004$ \\
$\hat\varphi_L$~($\tilde m_{t-1}$,长动量) & $+0.689$$^{**}$ & $0.013$ \\
$\hat\varphi_2$~($\tilde\sigma_{t-1}$,波动溢价) & $+0.461$ & $0.132$ \\
$\hat\varphi_3$~($\tilde\sigma_{t-1}\!\cdot\! r_{t-1}$,交互) & $+0.377$ & $0.227$ \\
$\hat\alpha+\hat\beta$(GARCH持续性) & $-0.434$ & $0.158$ \\
MZ $R^2$(波动可预测) & $+0.362$ & $0.248$ \\
VR(5)(方差比) & $+0.347$ & $0.270$ \\
买持夏普 & $+0.488$ & $0.108$ \\
\bottomrule\end{tabular}
\begin{tablenotes}\scriptsize\item 实测效果序数化:中证1000=4,中证500=3,弱/一般=2,商品=1;未实测资产不参与。
\end{tablenotes}\end{threeparttable}\end{table}


表~\ref{tab:ashare} 至表~\ref{tab:stk} 报告分组估计,表~\ref{tab:rank} 将各参数与本项目已实测的择时策略绩效(遗传规划因子库 + IC 加权,2018--2026 回测)做 Spearman 秩相关。四点结论:

\textbf{(1)~短期动量 $\varphi_1$ 是第一判别变量。}秩相关 $+0.77$($p=0.004$)。A股宽基呈单调梯度:中证1000($+0.055$,$t=2.7$)$>$ 中证500 $>$ 沪深300(不显著),与实测``爆炸$>$中等$>$弱''完全一致;商品 $\varphi_1$ 一致为负(日频反转,幅度 4--6bp,低于双边交易成本),美股为负且伴随高漂移。

\textbf{(2)~波动溢价 $\varphi_2$ 是第二来源。}中证1000 为 $+11.0$bp/日($t=2.8$),沪深300 为负。这一项直接对应因子库的构成:zz1000 在役因子中约 $60\%$ 为波动族(信号与滚动波动率高度相关且方向为``高波动做多''),其 cos IC 排名前 20 位全部属于此类。市场有什么结构,遗传搜索就长成什么形状。

\textbf{(3)~波动可预测性与择时能力无关。}GARCH 持续性与 MZ $R^2$ 的秩相关均不显著;贵州茅台与纳斯达克100 的波动最可预测(MZ $R^2\approx0.27$--$0.32$)而方向最不可测。``可以预测波动''不等于``可以择时''。

\textbf{(4)~个股反转与指数动量并存:聚合的涌现性质。}六只个股 $\varphi_1$ 全部为负(华鲁恒升 $-0.055$,$t=-2.7$,幅度与中证1000 相同而符号相反)。指数层面的正自相关并非成分股各自具有动量,而是特质反转在横截面上相互抵消后,共同成分的 lead--lag 传导所致(Lo--MacKinlay)。\emph{可择时性存在于聚合层,不可下沉至成分股。}

\section{机制与项目内证据的互证}
\textbf{因子库族谱。}按信号与回归特征的实际相关归类,zz1000 在役 79 个因子中:波动族 47 个(平均 cos IC $+0.079$)、动量族 17 个(平均 $-0.018$)、交互族 0 个。交互族缺位与 $\varphi_3$ 不显著($t=0.25$)一致——市场中不存在该结构,进化未产生对应个体,属正确的``未过拟合''。

\textbf{长动量与``元老''因子的方向锁死。}$\varphi_L$ 在中证1000 上为弱正($+5.4$bp,$t=1.5$)。项目中 9 个人工动量因子(``元老'')的方向由 2013--14 预热期一次性设定为 $-1$ 且从不更新,恰与 2015 年后的正向长动量相反:相同变换下原始动量信号 cos IC 为 $+0.051$,元老实测为 $-0.054$,互为镜像。剔除 7 个在役元老后组合年化自 $20.7\%$ 升至 $22.8\%$(夏普 $1.01\to1.11$);改用 cos IC 加权并排除元老为 $24.9\%$(夏普 $1.22$)。

\section{可择时性的操作化定义与推广}
\begin{quote}
\textbf{定义}(针对日频动量族多空择时算法):资产可择时,当且仅当
(T1) $\varphi_1>0$ 且 $t>2$、$\mathrm{VR}(5)>1.05$;
(T2) $\varphi_2>0$ 且 $t>1.5$;
(T3) 买入持有夏普 $|SR|<0.3$(低漂移,择时不缴``漂移税'')。
\end{quote}
23 个资产中同时满足三条的仅有:中证1000、中证500(各项弱一档)、\textbf{国证2000 与中证2000}。其中中证2000 各项均不弱于中证1000($\varphi_1=+0.054$,$\varphi_2=+10.8$bp,$\varphi_L=+10.6$bp 且显著,VR(5)$=1.205$ 为全样本最高),是最优先的推广标的;北证50 样本过短(853 日),暂不评价。需要提示:2000 系指数无对应股指期货,空头腿须以纯多口径或 IM(中证1000 期货)代理对冲实现。

\section{局限}
(i) 两步估计中状态变量为生成回归量,$\varphi_2,\varphi_3$ 的 $t$ 值存在轻度高估(Pagan, 1984);(ii) 秩相关样本仅 12 个资产,且实测等级为序数主观赋值;(iii) 黄金的择时表现集中于 2022 年以后,与本文``黄金四项皆不显著、仅有高漂移''的结论一致——该段收益应归因于趋势暴露而非择时技能;(iv) 回归均为全样本内估计,推广标的(2000 系)的结论尚待滚动样本外回测确认。

\end{document}
